\chapter{Apéndices}

\section{Código Fuente del Firmware (ESP32)}
Este código gestiona la lectura de todos los sensores (BME280, MPU6500, INMP441, Anemómetro) y envía los datos a Supabase mediante HTTPS.

\begin{lstlisting}[language=C++]
#include <WiFi.h>
#include <HTTPClient.h>
#include <Wire.h>
#include <Adafruit_Sensor.h>
#include <Adafruit_BME280.h>
#include <Adafruit_MPU6050.h>
#include <driver/i2s.h>
#include <time.h>
#include <ArduinoJson.h>

const char* ssid = "TU_SSID";
const char* password = "TU_PASSWORD";
const String supabase_url = "TU_URL";
const String supabase_anon_key = "TU_KEY";

void setup() {
  Serial.begin(115200);
  WiFi.begin(ssid, password);
  configTime(0, 0, "pool.ntp.org"); 
  Wire.begin(21, 22);
  bme.begin(0x76, &Wire);
  mpu.begin(0x69, &Wire);
  // ... (ver codigo completo en repositorio)
}

void loop() {
  // Lectura de microfono I2S y envio temporizado
  // ...
}
\end{lstlisting}

\section{Código Fuente del Frontend: Hook \texttt{useStationData}}
Este Hook de React centraliza la lógica de obtención de datos, manejo de históricos con paginación y suscripción en tiempo real.

\begin{lstlisting}[language=JavaScript]
import { useState, useEffect, useCallback, useRef } from 'react';
import { supabase } from '../lib/supabase';

export function useStationData(rangoInicial = '24H') {
  const [ultimaLectura, setUltimaLectura] = useState(null);
  const [historial, setHistorial] = useState([]);
  
  const fetchHistorial = useCallback(async (rangoActual) => {
    // Implementacion de paginacion para saltar limites de API
    let allData = [];
    let from = 0;
    const step = 1000;
    // ... logic loop
  }, []);

  useEffect(() => {
    const canal = supabase
      .channel('mediciones-realtime')
      .on('postgres_changes', { event: 'INSERT', schema: 'public', table: 'mediciones' }, 
      (payload) => {
        // Actualizacion reactiva del estado
      })
      .subscribe();
    return () => supabase.removeChannel(canal);
  }, []);

  return { ultimaLectura, historial, setRango };
}
\end{lstlisting}

\section{Esquema de Base de Datos (SQL)}
Definición de la tabla \texttt{mediciones} en PostgreSQL (Supabase).

\begin{lstlisting}[language=SQL]
CREATE TABLE mediciones (
  id uuid DEFAULT gen_random_uuid() PRIMARY KEY,
  created_at timestamptz DEFAULT now(),
  fecha_rtc timestamptz,
  temperatura_bme float4,
  humedad float4,
  presion float4,
  viento_pulsos int4,
  nivel_ruido float4,
  rssi_wifi int4,
  accel_x float4, accel_y float4, accel_z float4,
  gyro_x float4, gyro_y float4, gyro_z float4
);
\end{lstlisting}
